\documentclass[11pt]{article}

\usepackage[margin=1.05in,top=1.2in,bottom=1.15in]{geometry}
\usepackage{specimen}   % shared "Specimen Reclassified" design system, docs/specimen.sty

\setrunninghead{PROJECT OVERVIEW / KOSELLECK MACHINE}

\hypersetup{
  pdftitle={The Koselleck Machine: Project Overview},
  pdfsubject={What the project is, who it is for, where it stands, and how to run it},
}

\begin{document}

\specimencover
  {The Koselleck Machine}
  {A project overview: what this instrument measures, who it is for, what
   it can and cannot support yet, and how to run or extend it.}
  {%
  \coverfield{Repository}{koselleck-networks}
  \coverfield{Document}{docs/overview.tex $\rightarrow$ docs/overview.pdf}
  \coverfield{Companion}{docs/method.pdf (the method in full)}
  \coverfield{Status}{research instrument, local only, not yet deployed}
  \coverfield{Compiled}{\today}
  }

{\hypersetup{linkcolor=accentnet}
\tableofcontents}

\newpage

\section{What this is}

The Koselleck Machine is a research instrument built to test one
historical claim, and a small web tool for reading the result.

The claim is Reinhart Koselleck's. He argued that the decades between
roughly 1770 and 1830, which he called the \term{Sattelzeit} (German for
``saddle period''), were not merely a stretch of time in which many words
happened to change meaning at once, but a moment when the vocabulary of
politics and society reorganized \term{as a system}: concepts moved in
relation to one another rather than drifting independently. That has
always been argued in prose, from close reading. Ryan Heuser's
\emph{Computing Koselleck} was the first computational test of it: he
trained a separate word-embedding model per historical period and
measured how far each word moved from one period to the next, and found
that movement is unusually large around 1770-1830. But a
one-word-at-a-time measurement cannot distinguish a system reorganizing
from a pile of unrelated shifts that happen to coincide in time. This
project adds the missing test. It builds a word-similarity network for
each period, runs community detection on it, and measures whether the
\term{grouping structure itself} reorganizes at the Sattelzeit, rather
than only measuring individual words.

The full argument, the algorithms, and the exact measurements are set out
in the companion document, \code{docs/method.pdf}. This overview does not
repeat them. Its job is orientation: what the thing is, who it is for,
what state it is in, and how to run it.

\section{Who it is for, and what it is for right now}

Two readers are in scope.

\runin{A historian or intellectual historian} who wants to interrogate
the result rather than take a summary statistic on faith. For that
reader, the web tool is the product: type a word, see which semantic
community it belonged to in every twenty-year window from 1500 onward,
see explicitly where it moved and where it stayed, and drill into any
period to read the actual nearest neighbours behind that judgement.

\runin{A computational or digital-humanities reader} who wants to check
the method, disagree with it, or run it on different text. For that
reader, the product is the pipeline: seven standalone scripts, a single
versioned configuration file, and a corpus contract small enough that
substituting your own sources is a folder drop or one reader function
(Section~\ref{sec:byoc}).

\begin{claimbox}
\claimlabel{What it is not, yet.} This is a research instrument, not a
finished, citation-ready dataset. The numbers it produces are the object
under test, not a published result; the community labels are a reading
aid, not a validated taxonomy; and a full rerun does not reproduce an
identical partition. Section~\ref{sec:status} states each of those
plainly. Treat anything read out of the tool as a lead to check, not as
evidence to cite.
\end{claimbox}

\section{What is in the repository}

\runin{\code{src/}} is the pipeline, run in order: \code{parse\_tcp.py}
(raw sources to normalised per-document records), \code{bucket\_periods.py}
(records to 20-year text windows), \code{embeddings.py} (one word2vec
model per window), \code{network.py} (each word linked to its 15 nearest
neighbours by cosine similarity), \code{community.py} (Leiden community
detection at seven resolution settings), and \code{metrics.py} (how much
the grouping changed between consecutive periods: migration fraction,
NMI, adjusted Rand). Alongside them sit \code{pipeline\_config.py}, the
shared configuration loader, and the two optional labeling scripts
described below.

\runin{\code{webapp/}} is the Koselleck Machine itself, a Flask
application with a landing page and three views onto the same pre-built
networks: \code{/timeline} (the primary view: one word's community across
every period, in a single strip), \code{/graph} (a D3 force-directed
explorer for one word's neighbourhood in one period, plus a sampled
full-network mode), and \code{/search} (a plain lookup table: nearest
neighbours, community, and whether the word's community changed since the
previous period). Where the pipeline was also run per region, every page
gains a region toggle, listing only the regions that were actually built.

\runin{\code{docs/}} holds this overview and the method note.

\runin{\code{labels/}} holds a small, citable snapshot of the current
community labels, one CSV and one compiled JSON per region, copied there
by \code{label\_communities.py publish} so they travel with the code
instead of living only in the gitignored data directory.

\runin{\code{config.yml}} holds the settings that are shared and
versioned: the period slices, and the word2vec and Leiden
hyperparameters. Period boundaries and model parameters are configuration
here, not code.

What the repository deliberately does \emph{not} ship is the corpus or
the trained models. That is a size decision, not a rights one: the raw
Text Creation Partnership zips and the derived per-period networks run to
many gigabytes, and the sources are public domain and fetchable directly.
The contribution is the measurement, not the data.

\section{Where it stands}
\label{sec:status}

Three limitations bound what can be claimed today. All three are stated
in full in \code{docs/method.pdf} and surfaced inside the tool itself
rather than hidden; they are summarised here without softening.

\begin{claimbox}
\claimlabel{Corpus coverage.} The corpus is the Text Creation
Partnership: EEBO-TCP, ECCO-TCP, and Evans-TCP, curated and manually
corrected, with no OCR noise, but ending in 1800. The Sattelzeit runs to
1830. A Project Gutenberg supplement for 1800-1900 is planned and
\emph{not yet implemented}: there is no ingestion script for it, so the
pipeline as it stands covers 1500-1800 and later periods stay empty. The
gap is worst exactly where it matters most. It also skews composition and
not only volume: ECCO's public transcriptions thin out sharply after
1780, so the 1780-1800 window holds more American documents (Evans-TCP,
1681) than British ones (ECCO-TCP, 910). That is why the whole
measurement is also rerun per region, so a shift can be checked against
the possibility that it reflects a change in who was writing rather than
a change in the language. The per-region rerun bounds the problem rather
than solving it: each regional sub-corpus is thinner still than the
combined one, so its results carry wider uncertainty.
\end{claimbox}

\begin{claimbox}
\claimlabel{Label reliability.} Every community carries a short
plain-English name, and those names are generated by a language model
reading the community's most representative words. A label is inherited,
not re-read, for as long as the same alignment that decides
``moved''/``stayed'' maps a community back to a predecessor; only a
minority of labels (roughly one in ten to one in six, depending on
region) are an independent read in a given run. That independent read is
a single, unvalidated pass, useful for orientation and not for citation.
Inheritance carries a different cost: on the longest-running lineages a
label can go up to fifteen periods, close to three hundred years, without
being re-read, while the community's actual word composition drifts past
what the label still says. Separately, roughly 60-70\% of communities are
honestly flagged as ``mixed'', because the corpus is multilingual
(English, Latin, Law French, Welsh, Scots, and more) and early modern, and
a detected community sometimes groups words by shared grammar or shared
language rather than by shared topic. The tool shows that caveat next to
the labels rather than presenting them all as equally reliable.
\end{claimbox}

\begin{claimbox}
\claimlabel{Reproducibility.} Rerunning the pipeline from scratch on the
same corpus does not reproduce an identical result. Community detection
is properly seeded and deterministic given a fixed network; the gap sits
one step earlier, in training, which uses multiple parallel threads and
is therefore not bit-for-bit repeatable even with a fixed seed, because
thread scheduling varies. Small differences in the vectors cascade into
slightly different edges and then into a structurally similar but not
identical partition: one full rerun found 304 communities in the combined
corpus where an earlier run found 308. That is a real difference, not a
bug, and it is not worth trading away the threefold speed of parallel
training to chase, because the claim under test is whether reorganization
concentrates at the Sattelzeit and holds across the seven resolution
settings within a run, not whether a specific community keeps its id
across two independently retrained copies. The one thing designed to stay
comparable across reruns is the fixed lane list, because it is authored by
hand once and never reinvented.
\end{claimbox}

There is also a standing methodological rule worth repeating here,
because it governs everything the project will ever report:

\begin{claimbox}
\claimlabel{Hard rule.} Community detection resolution changes the
answer. Every network is therefore processed at seven resolution
settings, and no finding is reported if it holds only at a single,
convenient one. The sweep exists to test that a conclusion survives
across settings, not to pick the setting that agrees with the hypothesis.
\end{claimbox}

Finally, on availability: the Koselleck Machine currently runs locally
only. A public hosted deployment is intended but has not been started,
deliberately kept out of scope until the local application is solid. The
Flask app runs as-is on any host that can run Python, but it is not a
static site and cannot be served from GitHub Pages, since every response
is computed server-side from the network data.

\section{Running it}

\subsection{Setup}

\begin{dobox}
\code{python -m venv .venv}\\
\code{.venv/Scripts/activate}\quad\textrm{(or \code{source .venv/bin/activate} on macOS or Linux)}\\
\code{pip install -r requirements.txt}
\end{dobox}

The pipeline expects a \code{data\_root} folder containing
\code{corpus/}, \code{processed/}, \code{embeddings/}, \code{networks/},
and \code{communities/}. Point at your own copy either by creating a
gitignored \code{config.local.yml} with a \code{data\_root:} entry, or by
setting the \code{DATA\_ROOT} environment variable, which takes priority
over both config files. Only \code{corpus/} is filled in by hand;
everything else is produced by the pipeline. The exact folder names come
from \code{config.yml}'s \code{paths} block, which is authoritative if
anything is renamed.

The corpus itself is fetched separately. All three TCP components used
here are public domain, and the bulk P4 XML zip shards are downloaded
directly from the Text Creation Partnership. They are laid out under
\pathref{corpus/tcp/regions/<region>/<source>/*.zip}, for example
\code{british/eebo\_phase1/}, \code{british/ecco/},
\code{american/evans/}. Both folder levels are free-form:
\code{parse\_tcp.py} reads the zips in place, without unpacking them, and
discovers regions and sources by scanning that tree rather than from any
hardcoded list. The top-level folder name becomes the region tag used
throughout the pipeline and the webapp's region toggle. Only
quality-checked ``released'' shards are read; TCP's ``unedited'' variants
are skipped on purpose.

\subsection{Running the pipeline}

Each stage is a standalone script, run in order:

\begin{dobox}
\code{python src/parse\_tcp.py}\\
\code{python src/bucket\_periods.py}\\
\code{python src/embeddings.py}\\
\code{python src/network.py}\\
\code{python src/community.py}\\
\code{python src/metrics.py}
\end{dobox}

That sequence is enough to reproduce every quantitative result. Community
labels are a separate, optional stage on top of it, run as two more
scripts:

\begin{dobox}
\code{python src/extract\_community\_words.py}\\
\code{python src/label\_communities.py generate {\dd}region combined}\\
\code{python src/label\_communities.py compile {\dd}region combined}\\
\code{python src/label\_communities.py publish {\dd}region combined}
\end{dobox}

\code{generate} produces a human-editable CSV, inheriting each
community's label from its aligned predecessor wherever one exists and
leaving only genuinely new communities blank for a person or an agent to
fill in; \code{compile} turns the CSV into the JSON the webapp reads; and
\code{publish} copies the CSV and JSON into the repository's
\code{labels/} directory as a citable snapshot. \code{{\dd}region} also
accepts a region name or \code{all}. The prompt the model is given lives
in the repository as \pathref{src/prompts/community_labeling_v1.md}, and
\code{docs/method.pdf} quotes and documents it in full.

\subsection{Running the web tool locally}

\begin{dobox}
\code{python webapp/app.py}
\end{dobox}

Then open \url{http://127.0.0.1:5000}. Without real data the app still
starts, but every period shows as a coverage gap. For a production host,
the app is served with \code{gunicorn} and a \code{DATA\_ROOT}
environment variable pointing at wherever the network and community data
lives on that machine; the data does not travel with the repository and
has to be uploaded or fetched there separately.

\subsection{Bringing your own corpus}
\label{sec:byoc}

The pipeline is built so that a reader can rerun the whole measurement on
their own text rather than take these numbers on trust. There are exactly
two cases, and only one of them touches code.

\runin{Case A: the material is already in TCP's P4 XML format} (another
TCP release, or a repackaged TCP subset). No code change at all. Drop the
zip files at
\pathref{<data_root>/corpus/tcp/regions/<a-region-name>/<a-source-name>/*.zip}
and rerun the pipeline from \code{src/parse\_tcp.py}. Nothing needs
registering anywhere: both folder names are free-form, the region level
becomes the region tag carried through the pipeline and the webapp, the
source level is a diagnostic label kept in the manifest, and
\code{parse\_tcp.py} discovers both by scanning.

\runin{Case B: the material is in any other raw format} (plain text,
JSON, a different XML flavour, OCR dumps). One file changes:
\code{src/parse\_tcp.py}. Add a reader next to the existing
\code{iter\_xml\_members} that walks the new format and, for each
document, emits the same record the TCP path already writes to
\code{processed/all\_docs.jsonl}:

\begin{dobox}
\code{\{"region": "...", "source": "...", "doc\_id": "...", "year": 1782,
"text": "..."\}}
\end{dobox}

plus the matching \code{region,source,doc\_id,year,chars} row in
\code{manifest.csv}. Documents with no extractable year or no text are
dropped, as in the TCP path. That is the entire contract. Everything
downstream, \code{bucket\_periods.py} through \code{metrics.py} and the
webapp, reads only those normalised records and needs no change; period
boundaries and model hyperparameters are configuration in
\code{config.yml}, not code. In both cases, a region name that did not
exist before is picked up automatically by the region-split outputs and
the webapp's region toggle; there is no list of valid regions to update.

To be explicit about what does not exist: there is no upload interface,
no hosted ingestion service, and no API-key-driven pipeline. The
mechanism is the ordinary one, clone the repository, put your corpus
files in the right folder, rerun the scripts in order.

\section{People}

\runin{Ryan Heuser} wrote \emph{Computing Koselleck}, the word-level
antecedent this project extends, and is the author of the base method.

\runin{Jamie McGarry} pointed to the public TCP corpus and holds
Cambridge institutional access to the fuller restricted EEBO, ECCO, and
Evans archives beyond the public TCP subset, if a future window needs
more text than public TCP alone provides.

\runin{Bernardo Villegas} supervises the project and has given the
running feedback on method and on the design of the tool.

\section{Where to read more}

\code{docs/method.pdf} is the full method note: the argument, the
algorithms, the resolution sweep and why it is non-negotiable, the
labeling prompt and workflow in full, and the limitations in their
unabridged form. The repository's \code{README.md} carries the
operational detail: corpus download links and licensing, the exact
expected data layout, and deployment notes.

\end{document}
